\section{Digital Audio Restoration: Restoration of Global Degradations}
\subsection{Introduction}

The \textbf{restoration} of audio signals is important for several reasons.
Expectations on sound quality drive the need for accurate reproduction.
There is also a strong interest in preserving historical material, such as old records and old magnetic tapes, as well as reducing noise in modern digital recordings.

\begin{figure}[H]
    \centering
    \includegraphics[width=0.28\linewidth]{DAR 1.png}
    
\includegraphics[width=0.32\linewidth]{DAR 2.png}
    \includegraphics[width=0.3\linewidth]{DAR 3.png}
    \caption{Audio and media archives containing tapes, records and reels}
\end{figure}

Degradations in audio signals can be grouped into two main classes.
\textbf{Localized degradations} are discontinuities in the waveform that affect only certain samples (i.e. clicks, crackles, scratches, breakages or clipping).
\textbf{Global degradations} act on all samples of the waveform and include background noise, wow, flutter and certain types of non-linear distortion.

\subsection{Wiener filtering}
Regarding the \textbf{restoration of global degradations}, let us introduce a frequency-based perspective.
The \textbf{Wiener filter} is a signal processing technique used to recover a desired signal $s$ from an observed signal $x$ that has been corrupted by \textbf{additive noise} $v$: the signal $s$ and the noise $v$ are summed together, producing the input $x$ to a linear filter with impulse response $h$.  
The output of the filter, denoted $y = \hat{s}$, is an estimate of the original signal $s$.

\begin{figure}[H]
    \centering
    \includegraphics[width=0.5\linewidth]{WF degradation.png}
    \caption{Wiener filtering scheme for digital audio restoration}
\end{figure}

In this context, both the signal $s$ and the noise $v$ are assumed to be \textbf{stationary} and \textbf{uncorrelated} stochastic processes with zero mean.  
The noise $v$ is considered additive, meaning that it simply adds to the signal without affecting it in any multiplicative or non-linear way.

The main objective is to design the filter $h$ in such a way that the estimation error between the original signal $s$ and the filtered output $y$ is as small as possible in the mean squared sense.  
Formally, the problem consists of finding $h$ that minimizes the \textbf{Mean Squared Error (MSE)}, defined as:

\[
\mathrm{MSE} = E\left\{ (s - y)^2 \right\},
\qquad
E\{\cdot\}: \text{expectation operator}
\]

Minimizing this expression leads to the so-called \textbf{Minimum Mean Squared Error (MMSE)} condition, which provides the optimal Wiener filter solution.
The concept was first formalized by Norbert Wiener in 1949, who developed a complete mathematical framework for optimal filtering in the presence of noise.

\clearpage

\subsubsection{Orthogonality principle}

The MSE is minimized when the estimation error $e$ is \textbf{orthogonal} to the observed signal $x$:

\[
e = (s - y) \perp x
\]

In other words, the error signal is \textbf{uncorrelated} with the input signal used for estimation.  
Formally, this condition can be written as:

\[
E\left\{ \left[ s(n) - y(n) \right] \ x(m) \right\} = E\left\{ s(n) - y(n) \right\} \ E \left\{ x(m) \right\},
\qquad
\forall \ n, m
\]

If the signals involved are \textbf{zero-mean} processes, this simplifies to:

\[
E\left\{ \left[ s(n) - y(n) \right] \ x(m) \right\} = 0,
\qquad
\forall \ n, m
\]

This result is known as the \textbf{orthogonality principle} and it is a fundamental property of the Wiener filter: the optimal filter is such that its estimation error is orthogonal to the data used to obtain the estimate.

\subsubsection{Wiener filter response}

The orthogonality principle implies that:

\[
E\left\{ \left[ s(n) - y(n) \right] \ x(m) \right\} = 0
\quad \Rightarrow \quad
\underbrace{E\left\{ s(n) \ x(m) \right\}}_{R_{sx}} = \underbrace{E\left\{ y(n) \ x(m) \right\}}_{R_{yx}}
\]

That is, the \textbf{cross-correlation} $R_{sx}$ between the desired signal $s$ and the input $x$ must be equal to the \textbf{cross-correlation} $R_{yx}$ between the filter output $y$ and the input $x$.
Since the output $y$ is obtained by filtering the input $x$ through $h$, we can write:

\[
y(n) = \sum_k h(k)\ x(n-k)
\quad \Rightarrow \quad
R_{sx}(n-m) = E\left\{ \underbrace{\sum_k h(k) \ x(n-k)}_{y(n)} \ x(m) \right\}
\]

Because $h$ is \textbf{deterministic}, it comes out of the expectation, leading to:

\[
R_{sx}(n-m) = \sum_k h(k) \ E\left\{ x(n-k) \ x(m) \right\}
= \sum_k h(k) \ R_x(n-k-m),
\qquad R_x: \text{\textbf{auto-correlation} of } x
\]

Finally, by setting $m=0$ for simplicity:

\[
R_{sx}(n) = \sum_k h(k) \ R_x(n-k)
\]

Taking the Discrete-Time Fourier transform (DTFT) of both sides gives the frequency-domain Wiener filter:

\[
P_{sx} (\omega) = H(\omega) \ P_x (\omega)
\quad \Rightarrow \quad
H(\omega) = \frac{P_{sx}(\omega)}{P_x(\omega)}
\]

The term $P_{sx}(\omega)$ is the \textbf{cross-power spectral density} between $s$ and $x$, while $P_x(\omega)$ is the \textbf{power spectral density} of $x$.

\subsubsection{Considerations on power spectral densities}

As we assumed before, the desired signal $s$ and the noise $v$ are uncorrelated and both have zero mean.  
The \textbf{cross-correlation} $R_{sx}$ between original signal $s$ and noisy input $x$ is:

\[
R_{sx}(n) = E \left\{ s(k) \ x(k+n) \right\}
\]

The noisy input $x$ is the sum of clean signal $s$ and noise $v$:

\[
x(k+n) = s(k+n) + v(k+n)
\quad \Rightarrow \quad
R_{sx}(n) = E \left\{ s(k) \left[ s(k+n) + v(k+n) \right] \right\}
\]

By linearity of the expectation operator:

\[
R_{sx}(n) = E\left\{ s(k) \ s(k+n) \right\} + \cancel{E\left\{ s(k) \ v(k+n) \right\}}
\]

The second term is zero because $s$ and $v$ are uncorrelated and zero-mean:

\[
R_{sx}(n) = E\left\{ s(k) \ s(k+n) \right\} = R_s(n)
\quad \overset{\mathcal{F}}{\longrightarrow} \quad
P_{sx}(\omega) = P_s(\omega)
\]

This result simplifies the Wiener filter expression since the \textbf{cross-power spectrum} $P_{sx}$ reduces to the power spectrum $P_s$ of the clean signal $s$.

The \textbf{auto-correlation} $R_x$ of $x$ is:

\[
R_x(n) = E\left\{ x(k) \ x(k+n) \right\}
\]

The noisy input $x$ is the sum of clean signal and noise:

\[
x(n) = s(n) + v(n)
\quad \Rightarrow \quad
R_x(n) = E\left\{ \left[s(k) + v(k)\right] \left[s(k+n) + v(k+n)\right] \right\}
\]

Expanding the terms:

\[
R_x(n) = E\left\{ s(k) \ s(k+n) \right\} + \cancel{E\left\{ s(k) \ v(k+n) \right\}} + \cancel{E\left\{ v(k) \ s(k+n) \right\}} + E\left\{ v(k) \ v(k+n) \right\}
\]

Since $s$ and $v$ are uncorrelated and zero-mean, the cross terms vanish:

\[
R_x(n) = E\left\{ s(k) \ s(k+n) \right\} + E\left\{ v(k) \ v(k+n) \right\} = R_s(n) + R_v(n)
\quad \overset{\mathcal{F}}{\longrightarrow} \quad
P_x(\omega) = P_s(\omega) + P_v(\omega)
\]

Thus, the \textbf{power spectrum} $P_x$ of the noisy signal $x$ is the sum of the power spectra $P_s$ and $P_v$ of clean signal $s$ and noise $v$.
From all the previous derivations, we obtain the \textbf{optimal Wiener filter} in the frequency domain under the MMSE criterion:

\[
H(\omega) = \frac{P_{sx}(\omega)}{P_x(\omega)}
\quad \Rightarrow \quad
H(\omega) = \frac{P_s(\omega)}{P_s(\omega) + P_v(\omega)}
\]

This filter achieves MSE minimization (MMSE) between original signal $s$ and restored signal $\hat{s}$.  
Since $H(\omega)$ is real and positive, it acts as a \textbf{zero-phase low-pass filter}, attenuating frequency components where the noise dominates.
The quantities $P_s(\omega)$ and $P_v(\omega)$ denote the power spectral densities of the clean signal and the noise, respectively, which are assumed to be either known a priori or estimated from the data.

\subsection{Non-zero mean values}
So far, we assumed the signals involved in the Wiener filtering process to have \textbf{zero mean}: if the desired signal $s$ has a non-zero mean $m_s$, and the noise $v$ has a mean value $m_v$, we have to consider a correction in order to ensure proper operation of the filter.
In this case, the first step is to \textbf{subtract both mean values} $m_s$ and $m_v$ from the input $x$.

\begin{figure}[H]
    \centering
    \includegraphics[width=0.65\linewidth]{WF structure.png}
    \caption{Wiener filtering for digital audio restoration: non-zero mean generalization}
\end{figure}

The Wiener filter $H(\omega)$ is applied to this zero-mean version of the input.  
Once the filtering is completed, the \textbf{mean} value of the signal $m_s$ is finally \textbf{added back} to the output, so that the estimate $\hat{s}$ correctly reflects the mean value of the original signal.
This strategy guarantees that the Wiener filter works under the standard zero-mean assumption, while still being able to restore the correct mean at the output.

\subsection{Filter’s bandwidth and SNR impact}
In practice, the Wiener filter generally behaves as a low-pass filter.
This behaviour comes from its frequency response:

\[
H(\omega) = \frac{P_s(\omega)}{P_s(\omega) + P_v(\omega)}
\]

When the signal-to-noise ratio is high, the filter tends to pass the signal unchanged:

\[
P_s(\omega) \gg P_v(\omega)
\quad \Rightarrow \quad
H(\omega) \to 1
\]

On the other hand, when the signal-to-noise ratio is low, the filter strongly attenuates the input:

\[
P_s(\omega) \ll P_v(\omega)
\quad \Rightarrow \quad
H(\omega) \to 0
\]

This means that filtering has little effect when the SNR is already high, as the signal is clean enough.
On the other hand, when the SNR is very low, it is often better to discard the signal completely, since almost no useful information remains compared to the noise.

\begin{figure}[H]
    \centering
    \includegraphics[width=0.3\linewidth]{WF signal power.png}
    \includegraphics[width=0.3\linewidth]{WF noise power.png}
    \includegraphics[width=0.3\linewidth]{WF filter power.png}
    \includegraphics[width=0.3\linewidth]{WF filter power 2.png}
    \caption{Wiener filtering: power spectra $P$ of different signals and filter response $H$}
\end{figure}

\subsection{Power spectra characterization}
So far, we assumed that the \textbf{power spectra} $P_s$ and $P_v$ are known: in practice, this assumption is not very realistic.  
The \textbf{noise} spectrum $P_v$ is usually \textbf{easier to estimate}, since recordings often contain silent segments that can be used as references for the noise statistics.
By contrast, estimating the \textbf{signal} spectrum $P_s$ is \textbf{more challenging}.
In some situations, we can rely on clean signal samples and assume that they share the same power spectral density.
In other cases, $P_s$ may be inferred from portions of the recording where the noise is absent, whenever such regions exist.

A typical approach to estimate the power spectral density is through the \textbf{periodogram}.
Given $N$ frames $s_i(n)$ of the sound signal, we compute the Fourier transform $S_i(\omega)$ of each frame and estimate the PSD as the average of the squared magnitudes:  

\[
P_s(\omega) = \frac{1}{N} \sum_{i=1}^N \abs{S_i(\omega)}^2
\]

Another approach relies on \textbf{model-based estimation} techniques, where we approximate the spectral characteristics of the signal using parametric models such as autoregressive (AR), moving-average (MA), autoregressive moving-average (ARMA) or more sophisticated ARMAX models.

\subsubsection{Definition of Power Spectral Density (PSD)}
The \textbf{Power Spectral Density (PSD)} of a random process is usually defined as the Fourier transform of its autocorrelation function:

\[
P_s(\omega) = \mathcal{F}\{ R_s(n) \},
\qquad
R_s(n) = E \left\{s(k) \ s(k+n) \right\}
\]  

While this definition is rigorous, it is not always practical for direct computation.
A more useful formulation is provided by the \textbf{Wiener-Khinchin theorem}, which relates the PSD to the expected value of the squared Fourier transform of the signal:

\[
P_s(\omega) = E \left\{ \abs{S(\omega)}^2 \right\}
\]  

In practice, this expectation is approximated by computing the average over finite-length windows.
Considering $S_w(\omega)$ the Fourier transform of a windowed segment of the signal and $S_i(\omega)$ the Fourier transform of the $i$-th signal frame:

\[
P_s(\omega) = E \left\{ \abs{S(\omega)}^2 \right\} \approx E \left\{ \abs{S_w(\omega)}^2 \right\} 
\approx \frac{1}{N} \sum_{i=1}^N \abs{S_i(\omega)}^2
\]

This method of estimation is called the \textbf{periodogram} and it corresponds to averaging the squared magnitudes of the Fourier transforms of several overlapping or non-overlapping frames of the signal.

If the signal can be assumed \textbf{stationary} across the $N$ frames, the estimate becomes more reliable.
Otherwise, if stationarity cannot be guaranteed, the best approximation available reduces to the squared Fourier transform of a single windowed frame:  

\[
P_s(\omega) \approx \abs{S_w(\omega)}^2
\]

\clearpage

\subsection{Alternate formulation}
In practice, we do not have direct access to the clean signal $s$, but we can observe the noisy signal $x$ and the noise $v$.
For this reason, the Wiener filter can be reformulated in terms of the \textbf{available power spectral densities}.  
The output spectrum $\hat{S}(\omega)$ becomes:

\[
\hat{S}(\omega) = H(\omega) X(\omega) 
= \frac{P_s(\omega)}{P_s(\omega) + P_v(\omega)} X(\omega)
= \frac{P_x(\omega) - P_v(\omega)}{P_x(\omega) - \cancel{P_v(\omega)} + \cancel{P_v(\omega)}} X(\omega)
= \frac{P_x(\omega) - P_v(\omega)}{P_x(\omega)} X(\omega)
\]  

This shows that the filter response can be rewritten as follows:

\[
H(\omega) = \frac{\hat{S}(\omega)}{X(\omega)}
= \frac{P_x(\omega) - P_v(\omega)}{P_x(\omega)} 
= \frac{E \left\{ \abs{X(\omega)}^2 \right\} - E \left\{ \abs{V(\omega)}^2 \right\}}{E \left\{ \abs{X(\omega)}^2 \right\}} 
\approx \frac{\abs{X_w(\omega)}^2 - \abs{V_w(\omega)}^2}{\abs{X_w(\omega)}^2}
\]  

We approximate the expectations with finite-time estimates, where $X_w(\omega)$ and $V_w(\omega)$ are Fourier transforms of windowed segments of, respectively, noisy signal $x$ and noise $v$.

Notice that the numerator in this last expression may occasionally become negative due to random fluctuations.
Since the power spectral density of $s$ is by definition non-negative, such situations must be corrected, for instance by clipping negative values to zero.

\subsection{Implementation and performance}
The Wiener filter is typically implemented in the frequency domain, using the Discrete Fourier Transform (DFT) and its inverse (IDFT).
Since this implementation naturally performs a \textbf{circular convolution}, it corresponds to a true linear convolution only if the input sequences are adequately extended by \textbf{zero-padding}.
This prevents undesired wrap-around effects, which would otherwise distort the filtered output.

\begin{figure}[H]
    \centering
    \includegraphics[width=0.6\linewidth]{WF performance 1.png}
    \includegraphics[width=0.6\linewidth]{WF performance 2.png}
    \includegraphics[width=0.6\linewidth]{WF performance 3.png}
    \caption{Wiener filtering performance: original signal $s$ (top), noisy signal $x$ (center) and estimation $y$ (bottom)}
\end{figure}

In terms of \textbf{performance}, the Wiener filter is well known for its effectiveness in reducing noise.
It is theoretically optimal in the sense of minimizing the mean squared error (MMSE) between estimated and original signal.
However, one may question whether this criterion always aligns with perceptual relevance.

A limitation of the Wiener filter lies in its \textbf{stationarity}: it assumes that the signal and noise statistics remain constant over time.
This is not always true for real-world signals, which are often non-stationary.
As a result, the filter may introduce a \textbf{strong low-pass effect}, heavily impacting sound quality.
To mitigate these limitations, we adopt \textbf{adaptive versions} of the Wiener filter, which are able to track time-varying statistics and reduce artifacts.

\subsection{Power Spectral Matching (PSM)}
The Minimum Mean Square Error (MMSE) criterion is not the only \textbf{possible measure of optimality}, neither it is necessarily the most suitable one from a perceptual standpoint.
An alternative optimality criterion can be derived from the observation that the Wiener filter does not achieve exact power spectral matching.

\begin{figure}[H]
    \centering
    \includegraphics[width=0.5\linewidth]{WF degradation.png}
    \caption{Wiener filtering scheme for digital audio restoration}
\end{figure}

If the noisy signal $x$ is filtered by $H(\omega)$ to produce the estimate $\hat{s}$, the power spectral density of the output $y$ is:

\[
P_y(\omega) = P_x(\omega) H(\omega)^2 = P_x(\omega) \left( \frac{P_s(\omega)}{P_s(\omega) + P_v(\omega)} \right)^2
= \frac{\cancel{P_x(\omega)}}{\cancel{P_s(\omega) + P_v(\omega)}} \frac{P_s(\omega)^2}{P_s(\omega) + P_v(\omega)}
= \frac{P_s(\omega)^2}{P_s(\omega) + P_v(\omega)} \neq P_s(\omega)
\]

This shows that the Wiener filter output does not perfectly reproduce the spectral density $P_s$ of the original clean signal.
To enforce \textbf{Power Spectral Matching (PSM)}, we could instead define the filter response as follows:

\[
\hat{H}(\omega) = \sqrt{\frac{P_s(\omega)}{P_s(\omega) + P_v(\omega)}} 
\approx \sqrt{\frac{\abs{X_w(\omega)}^2 - \abs{V_w(\omega)}^2}{\abs{X_x(\omega)}^2}}
\quad \Rightarrow \quad
P_y(\omega) = P_x(\omega) \hat{H}(\omega)^2 = \frac{\cancel{P_x(\omega)} P_s(\omega)}{\cancel{P_s(\omega) + P_v(\omega)}} = P_s(\omega)
\]

The new filter definition $\hat{H}(\omega)$ guarantees that the spectral density of the restored signal matches that of the original source.  

\subsection{Parametric Wiener filtering}
The Wiener filter can also be expressed in the following equivalent form:

\[
H(\omega) = \frac{P_x(\omega) - P_v(\omega)}{P_x(\omega)} = 1 - \frac{P_v(\omega)}{P_x(\omega)}
\]

A useful extension of this concept is the \textbf{Parametric Wiener Filter}, which introduces two parameters: $\alpha$ and $\beta$. This formulation enables a \textbf{continuous transition} between the classical Wiener filter and the Power Spectral Matching (PSM) criterion.  
The filter is defined as follows:

\[
H_p(\omega) = \left[ \frac{P_x(\omega) - \alpha P_v(\omega)}{P_x(\omega)} \right]^{\beta}
= \left[ 1 - \alpha \ \frac{P_v(\omega)}{P_x(\omega)} \right]^{\beta}
\]

The parameter $\alpha$ controls the balance between noise suppression and signal preservation, while $\beta$ allows for shaping the aggressiveness of the filter response.
In practice, small values of $\beta$ lead to smoother transitions and less attenuation, whereas larger values of $\beta$ enhance suppression at low signal-to-noise ratios (SNRs).

\begin{figure}[H]
    \centering
    \includegraphics[width=0.33\linewidth]{WF parametric 1.png}
    \includegraphics[width=0.35\linewidth]{WF parametric 2.png}
    \caption{Parametric Wiener filtering: effect of $\beta$ (left) and $\alpha$ (right) on the filter response}
\end{figure}

The graphical representations confirm these effects: varying $\beta$ modifies the slope of the attenuation curve, while adjusting $\alpha$ determines the effective weight given to the noise component in relation to the signal.
This parametric formulation thus provides a flexible trade-off between the traditional Wiener filter and exact spectral matching.

\subsection{Background noise reduction}
In many practical scenarios, noise is modelled as an additive component, where the useful signal $s$ and the noise $v$ are assumed to be uncorrelated.
In this case, the observed signal $x$ and the Wiener filter output $y$ are:

\[
x = s + v,
\qquad
y = \hat{s} = H(\omega) \ x
\]

\textbf{Background noise}, such as the familiar hissing often heard in recordings, can be characterized as a single noise process that arises from several different sources.
These include electrical circuit noise, irregularities introduced by storage media, and ambient sounds present in the recording environment.

\begin{figure}[H]
    \centering
    \includegraphics[width=0.5\linewidth]{WF degradation.png}
    \caption{Wiener filtering scheme for digital audio restoration}
\end{figure}

Such background noise typically contains components across the entire frequency range.
In many cases, it can be modelled as a white random process: because of this, simple filtering or equalization methods are usually inadequate for effective suppression.
A more rigorous stationary solution is provided by the \textbf{Wiener filter}, which minimizes the mean square error between the clean and estimated signals.
However, when the noise characteristics vary over time, \textbf{adaptive strategies} are required.
A notable example is the \textbf{Short-Time Spectral Attenuation (STSA)} method, which operates in a time-frequency framework and adapts the filtering to local spectral conditions.

\subsubsection{Short-Time Spectral Attenuation (STSA)}
While stationary filtering methods, such as the Wiener filter, can be effective in some cases, they are usually inadequate when the noise statistics vary over time.
For this reason, an adaptive counterpart is often used, known as \textbf{Short-Time Spectral Attenuation (STSA)}.

\begin{figure}[H]
    \centering
    \includegraphics[width=0.65\linewidth]{STSA.png}
    \caption{Short-Time Spectral Attenuation (STSA) scheme}
\end{figure}

The method relies on a time-frequency analysis, typically based on the Short-Time Fourier Transform (STFT) or an equivalent filterbank.
The noisy signal $x$ is first decomposed into \textbf{frequency bands} using the STFT.
Within each band, the spectral components are attenuated according to a \textbf{suppression rule} that depends on the estimated noise power.
Finally, the processed bands are recombined through the inverse STFT to obtain the enhanced signal $y$.

The \textbf{attenuation function} $G(t_a^u, \Omega_k)$, at given time frame $t_a^u$ and frequency bin $\Omega_k$, depends on the power level $\abs{X(t_a^u, \Omega_k)}^2$ of the noisy signal measured at the same bin.
It also depends on the power $\hat{P}_v(\Omega_k)$ of the noise at frequency $\Omega_k$, assumed stationary and estimated as follows:

\[
\hat{P}_v(\Omega_k) = E \left\{ \abs{V(t_a^u, \Omega_k)}^2 \right\}
\]

The suppression rule ensures that the estimated clean spectrum is obtained by attenuating each frequency component in proportion to its estimated signal-to-noise ratio.

\subsubsection{Suppression rules}
Different suppression rules can be formulated depending on the chosen optimality criterion.  
The classical \textbf{Wiener suppression rule} is defined as follows:

\[
G(t_a^u, \Omega_k) = 1 - \frac{\hat{P}_v(\Omega_k)}{\abs{X(t_a^u, \Omega_k)}^2}
\]

The \textbf{PSM suppression rule} is defined as follows:

\[
G(t_a^u, \Omega_k) = \sqrt{1 - \frac{\hat{P}_v(\Omega_k)}{\abs{X(t_a^u, \Omega_k)}^2}}
\]

The \textbf{parametric Wiener suppression rule} is defined as follows:

\[
G(t_a^u, \Omega_k) = \left( 1 - \alpha \frac{\hat{P}_v(\Omega_k)}{\abs{X(t_a^u, \Omega_k)}^2} \right)^{\beta}
\]

We remind that $\alpha$ and $\beta$ are part of the parametric Wiener filter definition.

It is important to observe that in all of the above cases, the asymptotic behaviour of the suppression function is consistent:

\[
\lim_{\text{SNR} \to +\infty} G(t_a^u, \Omega_k) = 1,
\qquad
\lim_{\text{SNR} \to 0} G(t_a^u, \Omega_k) = 0
\]

\subsection{De-reverberation (channel inversion)}
In certain applications, such as speech analysis, \textbf{reverberations} are regarded as a problem since they act as a source of convolutive degradation.
Mathematically, the observed signal $x$ is the convolution of the dry speech signal $s$ with the \textbf{Room Impulse Response (RIR)} $h$:

\[
x(n) = h \ast s \ (n)
\]

The challenge is therefore to remove, or at least minimize, the impact of reverberation.  
The general algorithm can be summarized as follows:  

\begin{enumerate}
    \item Determine $h$: by direct measurement or by estimation from the reverberant speech $x$ (i.e. through cepstrum-based de-reverberation)
    \item Determine an \textit{inverse} representation of $h$
    \item Filter the observed signal $x$ with the inverse of $h$ in order to approximate the original clean signal
\end{enumerate}

To analyse de-reverberation in terms of channel inversion, we introduce the definition of the \textbf{complex cepstrum} $c_x$.

\begin{figure}[H]
    \centering
    \includegraphics[width=0.75\linewidth]{Cepstrum.png}
    \caption{Complex cepstrum definition}
\end{figure}

Starting from the observed signal $x$, we take the Fourier transform, apply the logarithm and compute the inverse Fourier transform:

\[
c_x(n) = \mathcal{F}^{-1}\left\{\log X(\omega)\right\} = \mathcal{F}^{-1}\left\{\log \left(H(\omega) \ S(\omega) \right)\right\}
\]

According to the logarithmic property:

\[
c_x(n) = \mathcal{F}^{-1}\left\{\log \left(H(\omega) \ S(\omega) \right)\right\}
= \mathcal{F}^{-1}\left\{\log \left(H(\omega)\right) + \log \left(S(\omega)\right)\right\} 
= \underbrace{\mathcal{F}^{-1}\left\{\log \left(H(\omega)\right)\right\}}_{c_h(n)} + \underbrace{\mathcal{F}^{-1}\left\{\log \left(S(\omega)\right)\right\}}_{c_s(n)}
\]

This shows that the complex cepstrum corresponds to a spectral analysis of the log-spectrum.
It separates the contribution of system $h$ and signal $s$, where low-quefrency terms typically represent the spectral envelope and high-quefrency terms represent fine spectral details.
These different spectral components can then be isolated and \textit{liftered} out in order to minimize the effect of reverberation.

In the log-spectrum representation, \textbf{speech components} tend to appear as \textbf{low-quefrency} contributions.
This means that the cepstral representation of the speech, denoted by $c_s$, is concentrated close to the cepstral origin.
By contrast, the \textbf{reverberant channel} is mainly represented by \textbf{high-quefrency} components.
The cepstrum of the channel, denoted by $c_h$, is made of pulses that are located farther away from the origin.
This separation between low and high-quefrency terms allows one to design de-reverberation schemes based on cepstral domain processing.

\begin{figure}[H]
    \centering
    \includegraphics[width=0.65\linewidth]{Dereverberation.png}
    \caption{De-reverberation scheme}
\end{figure}

The observed signal $x$ is first transformed into its \textbf{cepstrum}.
Once in the cepstral domain, we can inspect the coefficients to identify whether \textbf{peaks} related to \textbf{reverberation} are present.
If such peaks are detected, they can be suppressed by setting them to zero.
Alternatively, a \textbf{low-pass liftering} operation can be applied in order to preserve only the low-quefrency (speech-related) terms while attenuating the high-quefrency contributions (reverberation-related).
After the undesired components have been reduced or eliminated, the modified cepstrum is transformed back into the time domain, resulting in the de-reverberated signal $\hat{s}$.

\textbf{Channel inversion} refers to the concept of inverting the RIR $h$, obtaining a new filter $g$.
Let us consider the room impulse response $h$ to be a \textbf{Finite Impulse Response (FIR)} filter.
The goal is to estimate an inverse filter $g$:

\[
h \ast g \ (n) = \delta(n)
\]

In general, the inverse of an FIR filter is an Infinite Impulse Response (IIR) filter, which is not necessarily guaranteed to be stable.
To avoid stability issues, we typically seek an \textbf{approximate FIR} inverse by formulating the problem as a least-squares solution of an overdetermined system of linear equations:

\[
\underbrace{
\begin{bmatrix}
1 \\ 0 \\ \vdots \\ 0
\end{bmatrix}}_{\underline{d}}=
\underbrace{
\begin{bmatrix}
h(0) & 0 & \cdots & 0 \\
h(1) & h(0) & \ddots & \vdots \\
\vdots & h(1) & \ddots & 0 \\
h(N_h) & \vdots & \ddots & h(0) \\
0 & h(N_h) & \ddots & h(1) \\
\vdots & \ddots & \ddots & \vdots \\
0 & \cdots & 0 & h(N_h)
\end{bmatrix}}_{\left[ H \right]}
\underbrace{
\begin{bmatrix}
g(0) \\ g(1) \\ \vdots \\ g(N_g)
\end{bmatrix}}_{\underline{g}}
\]

The \textbf{least-squares estimate} of the inverse filter $g$ is computed as follows:

\[
\underline{g}_{LS} = \left[ \left[ H \right]^T \left[ H \right] \right]^{-1} \left[ H \right]^T \ \underline{d}
\]

Notice that if the impulse response $h$ is non-minimum phase (i.e. it has unstable zeros), then the residual error does not vanish even when the length of the inverse filter $N_g \to \infty$.

\clearpage
